\begin{figure}[!htbp]
\centering
\begin{tikzpicture}[>=Latex,font=\small,
  box/.style={draw=qs axis,rounded corners=2pt,align=center,text width=51mm,
             minimum height=10mm,inner sep=4pt,fill=white}]
  \node[box,text width=65mm] (density) at (0,0) {Präparierte Materialverschiebung\\\(\delta n_i\) und \(\delta T_{\mu\nu}\)};
  \node[box,fill=qs blue!7] (gravity) at (-3.55,-1.61)
     {Gewöhnliche Gravitation\\\(D_0\bar\gamma_{\mu\nu}=16\pi G\,\delta T_{\mu\nu}\)};
  \node[box,fill=qs rust!9] (higgs) at (3.55,-1.61)
     {Radiale Higgsquelle\\\(\delta\mathcal J_h\) aus W2};
  \node[box,fill=qs blue!7] (exterior) at (-3.55,-4.13)
     {Kompakter linearer\\mechanischer Beitrag};
  \node[box,fill=qs rust!9] (profile) at (3.55,-4.13)
     {Gerichtetes OP5-Profil\\\(\delta\mathcal J_{\rm OP5}=\mathcal P_{\rm dyn}\delta H\)};
  \node[box,fill=qs rust!9] (axis) at (3.55,-5.66)
     {Gezeitenkontrast auf dem Plateau\\\(\Delta\mathcal I\)};
  \draw[->,thick,qs blue] (density.south west) -- (gravity.north);
  \draw[->,thick,qs rust] (density.south east) -- (higgs.north);
  \draw[->,thick,qs blue] (gravity) -- (exterior);
  \draw[->,thick,dashed,qs rust] (higgs) -- (profile);
  \node[anchor=east,align=right,inner sep=4pt] at (3.52,-2.9)
      {angesetzte Quellenidentifikation\\[-1pt]\(\delta\mathcal J_{\rm OP5}=\delta\mathcal J_h\)};
  \draw[->,thick,qs rust] (profile) -- (axis);
\end{tikzpicture}
\caption{Die beiden Antwortwege derselben linearen Materialpräparation. Der gewöhnliche Beitrag folgt aus dem mechanischen Tensor. Der gerichtete Vergleich setzt die Kopplung der radialen W2-Quelle an das OP5-Profil voraus. Die gestrichelte Verbindung markiert diese Annahme.}
\label{fig:op5-modell}
\end{figure}
